%% Revised ICT Express manuscript. Build from this directory with latexmk -pdf.
\documentclass[a4paper,fleqn]{cas-dc}
\usepackage[numbers,sort&compress]{natbib}
\usepackage{booktabs}
\usepackage{graphicx}
\usepackage{amsmath}
\usepackage{url}
\usepackage{caption}
\usepackage{etoolbox}
\usepackage{needspace}
\hypersetup{pdftitle={LLM-Guided Bilingual RDF Retrieval: Alignment or Selection?},pdfauthor={Yuxiang Sun; Nan Yin},pdfsubject={Research article},pdfkeywords={knowledge graph, RDF retrieval, cross-lingual entity alignment, large language model}}
\setlength{\parskip}{0pt plus 1pt}
\setlength{\textfloatsep}{8pt plus 2pt minus 2pt}
\setlength{\intextsep}{8pt plus 2pt minus 2pt}
\newcommand{\Sstrict}{\ensuremath{S_{\mathrm{strict}}}}

\begin{document}
\let\WriteBookmarks\relax
\def\floatpagepagefraction{1}
\def\textpagefraction{.001}
\shorttitle{LLM-guided bilingual RDF retrieval}
\shortauthors{Y. Sun and N. Yin}
\title[mode = title]{LLM-Guided Bilingual RDF Retrieval: Alignment or Selection?}

\author[1]{Yuxiang Sun}[type=editor,orcid=0000-0003-0165-7664]
\credit{Conceptualization, Methodology, Software, Investigation, Writing -- original draft}
\ead{syx921120@gmail.com}
\affiliation[1]{organization={Guangzhou University of Software},city={Guangzhou},country={China}}
\author[2]{Nan Yin}[orcid=0009-0007-3506-9592]
\credit{Validation, Data curation, Writing -- review \& editing}
\ead{ny@caist.ca}
\affiliation[2]{organization={Canada Advanced Institute of Science and Technology},city={Ottawa},country={Canada}}
\cortext[1]{Corresponding author.}

\begin{abstract}
We examine how much retrieval value an LLM adds when prompts expose identifiers and candidate information. In a synthetic bilingual RDF benchmark and archived Wikidata outcomes, the reconstructed prompt contained the exact answer URI in 200/200 one-hop and 130/150 cross-lingual queries. Across 500 synthetic questions, a Qwen2.5-7B planner yielded 386 strict successes, compared with 350 for fixed-$k$ retrieval; the net gain was confined to the cross-lingual subset. In a stratified 60-query intervention, exact target-label matches were 36/40, 32/40, and 31/40 across full, URI-redacted, and semantic-only prompts in the URI-present stratum, and 1/20, 0/20, and 0/20 in the URI-absent stratum, where other cues could remain. On a synthetic ICT graph, Qwen and exact-code parsing selected 60/60 explicit asset codes; altered mappings reduced target and context retrieval to 30/60. Six archived Wikidata contrasts were non-significant after Holm correction; the held-out 19/21 result remains exploratory. These results measure constrained selection, not autonomous identity discovery or field performance.\end{abstract}

\begin{keywords}
knowledge graph \sep RDF retrieval \sep cross-lingual entity alignment \sep large language model \sep evaluation methodology
\end{keywords}
\maketitle
\hypersetup{pdftitle={LLM-Guided Bilingual RDF Retrieval: Alignment or Selection?},pdfauthor={Yuxiang Sun; Nan Yin},pdfsubject={Research article},pdfkeywords={knowledge graph, RDF retrieval, cross-lingual entity alignment, large language model}}

\section{Introduction}\label{sec:intro}
Language models are increasingly placed inside knowledge-graph retrieval pipelines to select entities, infer relations, or propose graph traversals. An end-to-end answer can be correct because the model resolved an entity, because the correct identifier was already among its candidates, or because a downstream executor supplied the answer. These mechanisms support different claims. In particular, selecting a URI from a prompt that already contains it is not evidence of autonomous cross-lingual identity discovery.

Research on language models with knowledge graphs, retrieval-augmented generation, and graph retrieval has examined complementary parts of this pipeline~\cite{pan2024unifying,ibrahim2024kgllm,lewis2020rag,edge2024graphrag,gao2024rag}. Cross-lingual entity alignment instead concerns identifying equivalent entities across language-specific graphs, using structural, embedding-based, or language-model methods~\cite{sun2018jape,zeng2021alignment,yang2024segalign}. Evidence-oriented evaluation work also shows why system outputs should be assessed against retrieved evidence rather than answer strings alone~\cite{thorne2018fever,beir2021}. For bilingual RDF retrieval, a central attribution question remains: under different URI-exposure and candidate-information constraints, what retrieval value does an LLM add over a deterministic process?

We examine this question in one synthetic bilingual RDF pipeline and a secondary re-analysis of saved Wikidata outputs. The synthetic graphs have heterogeneous schemas and disjoint URI namespaces, while planner prompts can contain truth-derived bilingual profiles and candidate anchors. We audit the complete application-level prompt message and distinguish explicit URI exposure from other identity cues. ICT network-operations knowledge graphs integrate heterogeneous sources for incident analysis~\cite{tailhardat2024noriaui}. If prompt-supplied URIs are mistaken for cross-language resolution, a stale or missing mapping in a device inventory, alarm graph, service catalog, or regional technical corpus can retrieve the wrong device/service, omit incident context, or misdirect troubleshooting. These are motivating failure modes, not measured outcomes. The Wikidata comparison is also bounded: the matched LLM receives the training-pool alignment map for cross-lingual queries, while the deterministic comparator applies a mapping derived from the same pool directly. The No explicit alignment table condition omits that explicit map but still receives Chinese query text and English candidate labels.

This paper makes three empirical contributions: a code-and-seed reconstruction of prompt URI exposure; paired re-analysis of synthetic and archived Wikidata outcomes; and a transparent account of which results can be regenerated versus reanalysed. We audit one system's information flow rather than propose a new retrieval or entity-alignment algorithm.

\section{Benchmark and evaluation}\label{sec:method}
\subsection{Synthetic benchmark and retrieval pipeline}
A deterministic generator builds two \texttt{rdflib} graphs~\cite{rdflib} containing 49,781 triples, 3,481 entities, and 8,218 labels or aliases. Graph A uses Chinese labels and one predicate schema; Graph B uses English labels and a different schema. Matched entities have different URIs, and no cross-graph identity predicate is included. The benchmark contains 500 questions: 200 one-hop attribute queries, 150 two-hop work-attribute queries, and 150 cross-lingual queries whose answer-bearing fact is in the other graph.

A label and alias index supplies candidate anchors to a Qwen2.5-7B planner served by Ollama at temperature zero, in JSON mode, with a 256-token output budget~\cite{qwen25}. The prompt serializes each candidate anchor with its URI, graph, label, and score. It also includes candidate profiles with canonical entity keys, Chinese and English surface forms, and URIs from both graphs, along with a predicate-schema summary containing local predicate names. Although the prompt section header refers to a name-to-description mapping, the reconstructed schema values are local names, not natural-language descriptions. Profiles are generated from gold entity records; the seeded memory coverage parameter is 0.92, with 3,183 of 3,481 entity profiles realized (91.44\%). These fields constitute privileged experimental information, not naturally acquired model knowledge. The answer stage deterministically renders retrieved triples and uses a decoy fallback. The experiment thus measures planning, candidate selection, and retrieval under this prompt construction, rather than open-ended answer generation.

The primary run used store seed 20260101, benchmark seed 20260102, and simulated-memory seed 20260104. Its retained log records 500 Ollama-backed full-system planner calls and no mock fallbacks. Three separate 200-query sensitivity runs used data seeds 20260101, 20260301, and 20260401. One reuses the primary data seed. Their distinct query-text overlaps with the primary run are 91, 0, and 1, respectively. They are sensitivity runs, not independent replications, and are not pooled with the primary sample.

\subsection{Metrics}
For query $q$, let $R(q)$ denote the ranked URI list, and let $\rho(q)$ be the rank of a gold URI, with $\rho(q)=\infty$ if none appears. Let $E(q)$ be the retrieved triples, $V^{*}(q)$ the generator-recorded gold values, $S(q)$ the triples logically scanned by the executor, and $A(q)$ the deterministic rendered answer. The rank metrics are
\begin{equation}
\begin{aligned}
\mathrm{Hits@1}&=\frac{1}{n}\sum_{q=1}^{n}\mathbf{1}[\rho(q)=1],\\
\mathrm{MRR}&=\frac{1}{n}\sum_{q=1}^{n}\frac{1}{\rho(q)}.
\end{aligned}
\label{eq:rank}
\end{equation}
where $1/\infty=0$. Let $V^{*}(q)$ contain the string forms of generator-recorded gold values for $q$. The object-grounded value-in-context indicator is
\begin{equation}
\begin{split}
h_{\mathrm{obj}}(q)=\mathbf{1}\!\big[&\exists(s,p,o)\in E(q):
\mathrm{str}(o)\in V^{*}(q)\\
&\land\ \mathrm{str}(o)\text{ is a substring of }A(q)\big].
\end{split}
\label{eq:obj}
\end{equation}
The implementation converts RDF objects with Python \texttt{str()} and compares lexical strings; it does not preserve or compare RDF datatype or language-tag annotations. A matching lexical value can occur in a triple with a different subject or predicate, and duplicate values are not disambiguated by triple identity. The query-level strict score and value-in-context (VIC) are
\begin{equation}
\begin{aligned}
\Sstrict(q)&=\mathbf{1}[\rho(q)=1]h_{\mathrm{obj}}(q),\\
\mathrm{StrictAcc}&=\frac{1}{n}\sum_{q=1}^{n}\Sstrict(q),\\
\mathrm{VIC}&=\frac{1}{n}\sum_{q=1}^{n}h_{\mathrm{obj}}(q).
\end{aligned}
\label{eq:strict}
\end{equation}
StrictAcc requires both a top-ranked gold URI and a gold lexical object value found in retrieved evidence and the deterministic rendered answer; VIC assesses only the object-value condition. This is string-based object-level grounding, not triple-level validation or full RDF semantic validation. Hits@1 and MRR measure URI ranking and do not imply grounding. No human or LLM judge is used for the synthetic strict score.

For the primary store of $|T|=49{,}781$ triples, the graph scan ratio (GSR) and its mean are
\begin{equation}
\begin{aligned}
\mathrm{GSR}(q)&=\frac{|S(q)|}{|T|},\\
\overline{\mathrm{GSR}}&=\frac{1}{n}\sum_{q=1}^{n}\mathrm{GSR}(q),\\
\overline{s}&=\frac{1}{n}\sum_{q=1}^{n}|S(q)|.
\end{aligned}
\label{eq:gsr}
\end{equation}
GSR records logical triples examined by the implementation; it is not physical I/O, memory use, cost, latency, or an asymptotic runtime estimate.

\subsection{Prompt audit and statistical analysis}
Let $P_{\mathrm{app}}(q)$ be the set of URI strings in the complete application-level user-message content assembled for query $q$, and let $u^{*}(q)$ be its gold URI. Explicit URI exposure is
\begin{equation}
z(q)=\mathbf{1}[u^{*}(q)\in P_{\mathrm{app}}(q)].
\label{eq:exposure}
\end{equation}
The prompt-construction code sends one user-role message and adds no application-level system message, few-shot examples, or other context. Its dynamic fields are the query, candidate-anchor JSON, candidate-profile JSON, and predicate-schema JSON; static instructions are also included. We reconstructed this complete application-level message from the saved seeds and code, rather than inspecting archived historical prompt snapshots, which were not retained. The audit checks all of its fields: URI-bearing strings can occur in the anchor and profile JSON; the reconstructed query text, local-name schema values, and static instructions contain no URI-form strings. We report $z(q)$ for one-hop and cross-lingual queries; two-hop work URIs are traversal targets rather than planner-candidate exposure targets. The Ollama runtime's internal chat template is outside this application-level reconstruction. Absence of the exact answer URI does not imply absence of identity information: an equivalent URI from the other graph, canonical keys, query semantics, labels, aliases, schema context, and model priors can remain informative. The audit measures explicit URI inclusion only, not these other cues.

For paired query outcomes, let $b$ and $c$ denote discordant pairs favoring systems A and B. The exact two-sided McNemar probability~\cite{mcnemar1947} is
\begin{equation}
p_{\mathrm{McN}}=\min\!\left(1,2\sum_{j=0}^{\min(b,c)}{b+c\choose j}2^{-(b+c)}\right).
\label{eq:mcnemar}
\end{equation}
We recalculated the six saved Wikidata contrasts and applied Holm's step-down correction across all six~\cite{holm1979}. For ordered raw probabilities $p_{(1)}\le\cdots\le p_{(6)}$, the adjusted value at rank $i$ is $p^{\mathrm{Holm}}_{(i)}=\min\{1,\max_{j\le i}(7-j)p_{(j)}\}$. The primary synthetic ablations retain their separate, prespecified three-comparison correction family. A zero-discordance result or a non-significant test does not demonstrate equivalence. We report two-sided 95\% Clopper--Pearson intervals for observed proportions $x/n$, with endpoints $[B^{-1}_{.025}(x,n-x+1),B^{-1}_{.975}(x+1,n-x)]$ and the standard boundary values for $x=0$ or $x=n$. Intervals are query-level and do not adjust for dependence from repeated entities or profiles. Sensitivity-run summaries are descriptive; three runs do not support broad uncertainty claims. The boundary between regenerable synthetic analyses and reanalysis-only Wikidata records is summarized in Table~\ref{tab:availability}.

We conducted a stratified, paired three-condition intervention on 60 cross-lingual queries (sample seed 20261120): 40 whose prompts contained the exact answer-bearing gold URI and all 20 in which that URI string was absent. Some latter prompts still contained the same entity's URI from the other graph or a canonical key, so this stratum is not identity-cue-free. The original response specification was replaced in all conditions with one shared instruction to output a candidate surface form, avoiding conflicting canonical-key and URI output demands; the full condition retained all original candidate information. URI redaction removed URI values from both candidate-anchor and candidate-profile fields while preserving canonical keys, bilingual surface forms, and schema. The semantic-only condition additionally removed canonical keys and anchor graph/score metadata, retaining the query, natural-language labels and aliases, and predicate schema. Qwen2.5-7B ran at temperature zero in all conditions. To keep scoring comparable when identifiers were withheld, we compared the model's selected entity surface form with the gold label in the opposite-language graph. This controlled intervention tests candidate selection under progressively reduced explicit identity cues; it is exploratory and does not establish autonomous entity discovery. Full prompts, sent prompts, completions, seeds, and query-level outcomes are archived.

We also generated a synthetic ICT mini-graph with 540 triples and 60 bilingual device, alarm, and service queries (20 per domain). Complete mappings yielded 60/60 exact target retrieval and context retention; deleting 15 mappings and redirecting 15 to another same-type entity reduced both to 30/60. We then ran Qwen2.5-7B at temperature zero in JSON mode (96-token output budget) on the same queries, supplying 20 bilingual candidates in the relevant domain. Each query contains its generated asset code, which the model was asked to select; a deterministic exact-code parser is the baseline. The selected code was passed through both mapping conditions. This probes candidate selection under an explicit identifier cue, not autonomous semantic discovery. The graph is synthetic and not an estimate of field error rates; scripts, prompts, completions, seeds, and query-level outputs are retained.

\section{Results}\label{sec:results}
\subsection{Primary benchmark and ablations}
The full planner achieved strict accuracy of 386/500 (0.772; 95\% exact CI 0.733--0.808) and Hits@1 of 404/500 (0.808). Strict accuracy was 181/200 (0.905) for one-hop, 140/150 (0.933) for two-hop, and 65/150 (0.433) for cross-lingual queries. Table~\ref{tab:main} presents the saved ablations and implementation-specific baselines. In a paired comparison, the full planner achieved 386/500 strict successes versus 350/500 for fixed-$k$ retrieval without an aligner; the discordant outcomes favored the planner on 65 queries and the baseline on 29 (exact two-sided McNemar $p=2.62\times10^{-4}$; Holm-adjusted $p=5.23\times10^{-4}$ within the prespecified three-ablation family). The baseline succeeded on all 350 combined one-hop and two-hop queries and on none of the 150 cross-lingual queries. Thus, the net overall difference is concentrated in cross-lingual queries, while the baseline is stronger on the one- and two-hop families; it is not a uniform advantage across query types. Removing predicate constraints reduced two-hop strict accuracy to zero. Removing localized pruning left the observed outcomes unchanged (386/500 strict successes, with no discordant pairs; $p=1$) while increasing GSR from 0.000365 to 1.0. This equality on one workload is not an equivalence result.

Across the three planned synthetic ablations, discordant counts $(b,c)$ were $(65,29)$ for fixed-$k$ retrieval without an aligner, $(140,14)$ for removal of predicate constraints, and $(0,0)$ for removal of pruning. Exact McNemar probabilities were $2.62\times10^{-4}$, $2.56\times10^{-27}$, and 1; Holm-adjusted values were $5.23\times10^{-4}$, $7.67\times10^{-27}$, and 1. The vector comparator uses deterministic hashed character n-gram features. The label-lookup comparator is a pattern-matching implementation, not a general SPARQL~1.1 engine.

\begin{table*}[t]
\centering\small
\begin{tabular}{lcccc}
\toprule
System & Strict accuracy (95\% exact CI) & Hits@1 & GSR & Scanned triples/query \\
\midrule
LLM planner & 386/500 (0.772; 0.733--0.808) & 404/500 (0.808) & 0.000365 & 18.2 \\
No aligner (fixed $k=6$) & 350/500 (0.700; 0.658--0.740) & 350/500 (0.700) & 0.000298 & 14.8 \\
No predicate constraint & 260/500 (0.520; 0.475--0.565) & 260/500 (0.520) & 0.000306 & 15.2 \\
Global scan (no pruning) & 386/500 (0.772; 0.733--0.808) & 404/500 (0.808) & 1.000000 & 49,781 \\
Hashed-vector GraphRAG & 263/500 (0.526; 0.481--0.571) & 263/500 (0.526) & 0.000905 & 45.0 \\
Label-lookup pattern baseline & 200/500 (0.400; 0.357--0.444) & 200/500 (0.400) & 0.156118 & 7,772 \\
\bottomrule
\end{tabular}

\caption{Primary 500-query results. Strict accuracy is the object-grounded measure in Eq.~(\ref{eq:strict}) with exact 95\% intervals. GSR is logical triple access. The vector and label-lookup rows are implementation-specific comparators.}
\label{tab:main}
\end{table*}

\begin{table*}[t]
\centering\small
\begin{tabular}{lcc}
\toprule
Query family & LLM planner, strict & Fixed-k no-aligner, strict \\
\midrule
1-hop & 181/200 (0.905) & 200/200 (1.000) \\
2-hop & 140/150 (0.933) & 150/150 (1.000) \\
cross-lingual & 65/150 (0.433) & 0/150 (0.000) \\
\bottomrule
\end{tabular}

\caption{Descriptive strict-success rates for the full planner and fixed-$k$ no-aligner baseline by query family, calculated on the same query identifiers. No subgroup significance tests are reported.}
\label{tab:family}
\end{table*}

\subsection{Explicit URI exposure and candidate selection}
The corrected prompt reconstruction finds the exact answer-bearing gold URI in all 200 one-hop prompts (100\%) and in 130/150 cross-lingual prompts (86.67\%). On one-hop queries, the planner ranked the gold URI first on 195/200 and achieved strict accuracy on 181/200. Thus, this one-hop condition is not a gold-URI-unexposed retrieval test. For cross-lingual queries containing the exact answer-bearing URI, top-rank selection and strict accuracy were both 65/130 (50.0\%; 95\% exact CI for strict accuracy 0.411--0.589). Among the 20 prompts without that exact URI string, there were no top-rank or strict successes (0/20; 95\% exact CI for strict accuracy 0--0.168). Some of these prompts still contained a counterpart URI or canonical key, so this is not an identity-cue-free condition. In the synthetic prompt, profile mappings derive from truth records; prompt fields other than the exact answer URI can also reveal the target identity. Table~\ref{tab:probe} reports the corrected strata.

These measurements separate explicit URI exposure from other cues. Bilingual labels, aliases, query wording, and model priors may support inference, but were not isolated experimentally.

\begin{table*}[t]
\centering\small
\begin{tabular}{llccc}
\toprule
Family & Gold URI condition & $n$ & Strict accuracy (95\% exact CI) & Hits@1 \\
\midrule
1-hop & Gold URI present & 200 & 181/200 (0.905; 0.856--0.942) & 195/200 (0.975) \\
2-hop & Not applicable & 150 & 140/150 (0.933) & 144/150 (0.960) \\
cross-lingual & Gold URI present & 130 & 65/130 (0.500; 0.411--0.589) & 65/130 (0.500) \\
cross-lingual & No explicit gold URI & 20 & 0/20 (0.000; 0.000--0.168) & 0/20 (0.000) \\
\bottomrule
\end{tabular}

\caption{Primary results stratified by the exact answer-bearing URI string in the reconstructed prompt. URI absence does not imply absence of counterpart identifiers or other identity cues. Intervals are two-sided 95\% exact binomial intervals for strict accuracy. Two-hop work URIs are traversal targets, so exposure is not defined for that family.}
\label{tab:probe}
\end{table*}

\subsection{Three-condition URI intervention}
In the 40-query stratum containing the exact answer URI, exact opposite-language label matches were 36/40 with the full candidate data, 32/40 after removing all candidate URIs, and 31/40 after additionally removing canonical keys. In the 20-query stratum without that exact URI string, the corresponding counts were 1/20, 0/20, and 0/20. The latter stratum could still expose an equivalent URI or canonical key in the original prompt. The semantic-only condition also retained truth-derived bilingual candidate labels, so it measures label selection within a supplied profile, not autonomous identity discovery. These descriptive results use exact normalized surface-form matching; they are not evidence that masking improves inference or that conditions are equivalent. All 180 prompts and completions are archived.

\subsection{Synthetic ICT mini-graph probe}
Both selected all 60 exposed codes. Complete mappings yielded 60/60 correct targets and contexts; removing 15 mappings and redirecting 15 reduced both selectors to 30/60, with context loss in 30/60. Because queries exposed asset codes and Qwen received candidates, this synthetic probe measures constrained selection, not semantic discovery or field error.

\subsection{Sensitivity runs}
Across three 200-query sensitivity runs, strict accuracy ranged from 0.760 to 0.805 (mean 0.780; SD 0.023). Query-text overlaps with the primary run were 91, 0, and 1; one run reused its data seed, so these are not independent replications (Table~\ref{tab:seeds}).

\begin{table*}[t]
\centering\scriptsize
\begin{tabular}{l p{4.5cm}ccc}
\toprule
Data seed & Distinct query texts overlapping primary & Strict accuracy & Hits@1 & Value-in-context \\
\midrule
20260101 (primary seed) & 91 & 155/200 (0.775) & 0.810 & 0.915 \\
20260301 & 0 & 161/200 (0.805) & 0.830 & 0.955 \\
20260401 & 1 & 152/200 (0.760) & 0.825 & 0.915 \\
Mean (SD) & -- & 0.780 (0.023) & 0.822 & 0.928 \\
\bottomrule
\end{tabular}

\caption{Three sensitivity runs. Query overlap counts are distinct query texts shared with the primary run. The 20260101 sensitivity run reuses the primary data seed; the runs are not independent replications.}
\label{tab:seeds}
\end{table*}

\subsection{Archived Wikidata comparison}
The secondary archive contains 184 Chinese/English queries over 70 Wikidata entities, including 114 one-hop and 70 cross-lingual queries~\cite{wikidata2014}. The matched LLM receives the complete training-pool Chinese-label-to-URI map on cross-lingual queries; the deterministic system uses a mapping derived from the same training pool but applies a direct lookup rather than receiving the same serialized prompt. All paired systems use identical query identifiers and candidate pools. The No explicit alignment table condition receives query text and candidate labels without an explicit alignment table. For cross-lingual queries, these inputs still provide Chinese query semantics and English candidates; the condition therefore removes the explicit mapping but not all cross-language information.

The matched LLM and deterministic systems scored 157/184 and 163/184, respectively (Holm-adjusted $p=0.281$). On one-hop queries they scored 107/114 and 114/114 (adjusted $p=0.088$); on all cross-lingual queries, 50/70 and 49/70 (adjusted $p=1.000$). None of the six contrasts reached $p<0.05$ after the shared six-test Holm correction, including training and held-out subsets and the No explicit alignment table condition versus deterministic cross-lingual comparison (Table~\ref{tab:wikidata}). This lack of statistical significance does not show that the systems are equivalent.

Within the 21-query held-out cross-lingual subset, the No explicit alignment table condition scored 19/21, compared with 1/21 for the matched LLM and 0/21 for the deterministic system. We treat 19/21 as an exploratory observation only. The saved records and code do not establish whether held-out queries influenced prompt development, parameter choices, manual inspection, or experiment selection; successful raw model completions and exact prompt snapshots were also not retained. This result does not demonstrate autonomous cross-lingual entity identity discovery.

\begin{table*}[t]
\centering\scriptsize
\begin{tabular}{p{5.0cm}ccccc}
\toprule
Paired contrast & A strict successes & B strict successes & A-only/B-only & Exact $p$ & Holm-adjusted $p$ \\
\midrule
Matched LLM vs deterministic: all queries & 157/184 & 163/184 & 1/7 & 0.0703 & 0.281 \\
Matched LLM vs deterministic: one-hop & 107/114 & 114/114 & 0/7 & 0.0156 & 0.088 \\
Matched LLM vs deterministic: cross-lingual & 50/70 & 49/70 & 1/0 & 1 & 1.000 \\
Matched LLM vs deterministic: cross train & 49/49 & 49/49 & 0/0 & 1 & 1.000 \\
Matched LLM vs deterministic: cross held-out & 1/21 & 0/21 & 1/0 & 1 & 1.000 \\
\shortstack[l]{No explicit alignment table condition:\\LLM vs deterministic, cross-lingual} & 62/70 & 49/70 & 19/6 & 0.0146 & 0.088 \\
\bottomrule
\end{tabular}

\caption{Six paired contrasts recalculated from the archived query-level Wikidata outcomes. Counts are strict successes under the archived Wikidata evaluator. All Holm-adjusted probabilities use the same six-test family; none is below 0.05. The held-out 19/21 result is descriptive and exploratory, not evidence of autonomous alignment.}
\label{tab:wikidata}
\end{table*}

\section{Discussion and limitations}\label{sec:discussion}
In this synthetic pipeline, the planner's aggregate gain over fixed-$k$ retrieval comes from cross-lingual queries: the baseline succeeds on all one- and two-hop queries and none of the cross-lingual queries, while the planner succeeds on 65 cross-lingual queries. Exact answer URIs occur in 130/150 cross-lingual prompts; the remainder may still contain counterpart URIs or keys. The results concern retrieval with truth-derived profiles, not a broad independent LLM contribution. The archived Wikidata contrasts are non-significant; equivalence is unestablished.

Candidate selection chooses among supplied identifiers; semantic inference uses labels, relations, or schema to rank candidates. Autonomous discovery requires establishing cross-graph equivalence without direct mappings and with leakage and development-set reuse controlled. The semantic-only arm still received truth-derived bilingual labels, so the intervention tests label selection under reduced identifier cues. The mini ICT probe isolates missing or stale links in a synthetic graph, not LLM capability or field error.

The benchmark uses truth-derived profiles and deterministic answers; its lexical grounding ignores RDF language/datatype tags and evidence-triple subjects/predicates. One 7B model was tested, and GSR counts logical scans only. Sensitivity runs are not independent. Four Wikidata source inputs and historical prompts/completions are missing; held-out isolation is unverified. The new probes retain full prompts, outputs, inputs, seeds, and scripts (Table~\ref{tab:availability}); the ICT graph is synthetic.

\begin{table*}[t]
\centering\scriptsize
\begin{tabular}{>{\raggedright\arraybackslash}p{2.55cm}>{\raggedright\arraybackslash}p{12.15cm}}
\toprule
Material & What the retained archive supports \\
\midrule
Synthetic runs & Generator, seeds, configurations, outputs, and primary log support regeneration. Historical prompts/completions are absent; sensitivity logs are absent, overlaps are 91/0/1, and one run reuses the primary data seed. \\
Wikidata archive & Parsed outcomes, aggregates, manifest, and two LLM logs support reanalysis. Four source inputs and successful prompt/completion snapshots are missing, precluding source-level reruns. \\
New probes & URI probe retains 180 prompts/completions, strata, seeds, script, and outcomes; reruns require Ollama. ICT probe retains its graph, 60 Qwen prompts/completions, mapping outcomes, script, and seed; data are synthetic. \\
\bottomrule
\end{tabular}
\caption{Historical evidence limits and newly retained experimental records.}
\label{tab:availability}
\end{table*}

No operator or IoT workload is included. In ICT operations, a stale cross-source map could attach an alarm to the wrong asset, distort service-impact analysis, or omit troubleshooting context. Prior ICT knowledge-graph systems correlate heterogeneous sources~\cite{xie2021iotkg,tailhardat2024noriaui}; these are motivating failure modes, not observed deployment effects. Field validation requires held-out inventories, alarms, service catalogs, and multilingual technical documents.

\section{Conclusion}\label{sec:conclusion}
Planner gains over fixed-$k$ were confined to cross-lingual queries, 130/150 of which exposed the exact URI. The URI intervention yielded 36/40 to 31/40 exact-label matches when the URI was present and at most 1/20 when absent, though other cues remained. On the synthetic ICT graph, Qwen matched the exact-code parser (60/60); mapping corruption reduced retrieval and context retention to 30/60. Wikidata contrasts were non-significant; 19/21 is exploratory. Evaluation should audit identifiers, compare matched baselines, test mapping integrity, and track context retention. Autonomous discovery and field reliability remain unestablished.

\footnotesize
\section*{CRediT authorship contribution statement}
\textbf{Yuxiang Sun:} Conceptualization, Methodology, Software, Investigation, Writing -- original draft. \textbf{Nan Yin:} Validation, Data curation, Writing -- review \& editing.

\section*{Declaration of competing interest}
The authors declare that they have no known competing financial interests or personal relationships that could have appeared to influence the work reported in this paper.

\section*{Funding}
Funding: 2026 Research Project, Guangzhou University of Software (grant KY202649).

\Needspace{15\baselineskip}
\section*{Code and data availability}
Code, synthetic data, prompts, outputs, and scripts are archived in repository release v2.8. Wikidata records support reanalysis only; four source inputs and historical successful prompts are absent. Both new probes retain complete runs; rerunning the Qwen experiments requires Ollama. The ICT graph is regenerated from its script and seed.

\balance
\begin{thebibliography}{00}\footnotesize
\bibitem{pan2024unifying}
S. Pan, L. Luo, Y. Wang, C. Chen, J. Wang, X. Wu, Unifying large language
models and knowledge graphs: a roadmap, IEEE Trans. Knowl. Data Eng. 36 (7)
(2024) 3580--3599. https://doi.org/10.1109/TKDE.2024.3352100.
\bibitem{edge2024graphrag}
D. Edge, H. Trinh, N. Cheng, et al., From local to global: a graph RAG approach
to query-focused summarization, arXiv:2404.16130, 2024.
\bibitem{sun2018jape}
Z. Sun, W. Hu, Q. Zhang, Y. Qu, Bootstrapping entity alignment with knowledge
graph embedding, in: Proc. IJCAI, 2018, pp. 4396--4402.
https://doi.org/10.24963/ijcai.2018/611.
\bibitem{yang2024segalign}
L. Yang, J. Cheng, F. Zhang, Advancing cross-lingual entity alignment with
large language models, in: Findings of EMNLP 2024, ACL, 2024, pp. 8122--8138.
https://doi.org/10.18653/v1/2024.findings-emnlp.475.
\bibitem{thorne2018fever}
J. Thorne, A. Vlachos, C. Christodoulopoulos, A. Mittal, FEVER: a large-scale
dataset for fact extraction and VERification, in: Proc. NAACL-HLT, 2018,
pp. 809--819. https://doi.org/10.18653/v1/N18-1074.
\bibitem{lewis2020rag}
P. Lewis, E. Perez, A. Piktus, et al., Retrieval-augmented generation for
knowledge-intensive NLP tasks, in: Advances in Neural Information Processing
Systems 33 (NeurIPS), 2020, pp. 9459--9474.
\bibitem{rdflib}
D. Krech, E. Chuc, N. Car, A. Nelson, G. Higgins, A. Sommer, RDFLib,
version 7.6.0, 2026. https://doi.org/10.5281/zenodo.6845245.
\bibitem{qwen25}
A. Yang, et al., Qwen2.5 technical report, arXiv:2412.15115, 2024.
https://arxiv.org/abs/2412.15115.
\bibitem{mcnemar1947}
Q. McNemar, Note on the sampling error of the difference between correlated
proportions or percentages, Psychometrika 12 (2) (1947) 153--157.
\bibitem{holm1979}
S. Holm, A simple sequentially rejective multiple test procedure, Scand. J.
Stat. 6 (2) (1979) 65--70.
\bibitem{gao2024rag}
Y. Gao, Y. Xiong, X. Gao, K. Jia, J. Pan, et al., Retrieval-augmented
generation for large language models: a survey, arXiv:2312.10997, 2023.
https://doi.org/10.48550/arXiv.2312.10997.
\bibitem{zeng2021alignment}
K. Zeng, C. Li, L. Hou, J. Li, L. Feng, A comprehensive survey of entity
alignment for knowledge graphs, AI Open 2 (2021) 1--13.
https://doi.org/10.1016/j.aiopen.2021.02.002.
\bibitem{wikidata2014}
D. Vrande\v{c}i\'{c}, M. Kr\"{o}tzsch, Wikidata: a free collaborative
knowledgebase, Commun. ACM 57 (10) (2014) 78--85.
https://doi.org/10.1145/2629489.
\bibitem{beir2021}
N. Thakur, N. Reimers, A. R\"{u}ckl\'{e}, A. Srivastava, I. Gurevych, BEIR: a
heterogeneous benchmark for zero-shot evaluation of information retrieval
models, in: NeurIPS Datasets and Benchmarks Track, 2021. arXiv:2104.08663.
\bibitem{xie2021iotkg}
C. Xie, B. Yu, Z. Zeng, Y. Yang, Q. Liu, Multilayer Internet-of-Things
middleware based on knowledge graph, IEEE Internet Things J. 8 (4) (2021)
2635--2648. https://doi.org/10.1109/JIOT.2020.3019707.
\bibitem{tailhardat2024noriaui}
L. Tailhardat, Y. Chabot, A. Py, P. Guillemette, NORIA UI: efficient incident
management on large-scale ICT systems represented as knowledge graphs, in:
Proc. ARES 2024, GRASEC workshop, 2024.
https://doi.org/10.1145/3664476.3670438.
\bibitem{ibrahim2024kgllm}
N. Ibrahim et al., A survey on augmenting knowledge graphs (KGs) with large
language models (LLMs): models, evaluation metrics, benchmarks, and challenges,
Discov. Artif. Intell. 4 (2024) 76. https://doi.org/10.1007/s44163-024-00175-8.
\end{thebibliography}
\end{document}
